\section{The mathematics this problem still needs}

Everything in Parts I--IV of this document is either a theorem about
$\varphi$ or a theorem about proof methods for $\varphi$. Nothing in it
decides Erd\H{o}s \#249, and this section does not either. What this section
does is different in kind from the rest: it takes each route that survives
the barrier classification, states the missing mathematics as an exact
sentence, identifies the shape of argument that could supply it, names the
precise point at which the nearest existing technique fails, and --- where
the evidence permits --- attempts the construction rather than describing it.

Four things are new here and are marked as such. (i) An exact normal form
that turns every surviving \#249 route into a statement about the binary
expansion of one explicit real number, with the thresholds computed
(Lemma~\ref{lem:orbit}, Proposition~\ref{prop:transfer},
Corollary~\ref{cor:digitform}). (ii) A proof that the first-harmonic gap is
strictly stronger than irrationality, by an explicit witness inside the same
coefficient class (Theorem~\ref{thm:lacunary}); the corpus had this only as
a plausibility argument. (iii) A one-sided smooth-number majorant for the
non-supplier budget and a bad-cofactor estimate, with their remaining
uniformity hypotheses exposed (Propositions~\ref{prop:dickman} and
\ref{prop:badcof}); the quantifier audit also shows that the proposed
shallow-modulus Siegel--Walfisz route does not fit the predicate. (iv) An
audit of the rationality-side rank route: its proposed generic counterexample
loses one zero-residue section per level, so it neither refutes nor retires
the route.

Where a claim is proved it is marked \ev{Math}; where it is checked by exact
or floating-point computation, \ev{Cert}; where it is quoted, \ev{Cited};
where it is a proposal, \ev{Open}. No proposal below is offered as progress
towards a solution. Two of them are reductions, and a reduction is not a
result.

\begin{rem}
Declarations in the shared tree are cited as
\texttt{Erdos249257.name} with file and line. Declarations in the newer
per-problem tree are cited with their path from the repository root,
\texttt{ErdosProblems/Erdos249/File.lean:line}; the hyperlink target for
those is the per-problem directory, not \texttt{Erdos249257/}.
\end{rem}

\subsection{One identity, and what it does to every surviving route}

Write $S=\sum_{n\ge 1}\varphi(n)/2^n$, $R_N=\sum_{m\ge 1}\varphi(N+m)/2^m$
for the local tail
(\lean{Erdos249257.\allowbreak TotientTailPeriodKiller.\allowbreak totientTail}{TotientTailPeriodKiller.lean:57}),
and $\Phi_N=\sum_{n\le N}\varphi(n)2^{N-n}\in\Z$ for the integer prefix, so
that $2^N S=\Phi_N+R_N$. For $h\ge 1$ set
\[
  \alpha_h \;:=\; (2^h-1)\,S .
\]

\begin{lem}[Doubling normal form]\label{lem:orbit}
For all $N\ge 0$ and $h\ge 1$,
\[
  R_{N+1}=2R_N-\varphi(N+1),
  \qquad
  R_{N+h}-R_N \;=\; 2^N\alpha_h-\bigl(\Phi_{N+h}-\Phi_N\bigr),
\]
so $R_{N+h}-R_N \equiv 2^N\alpha_h \pmod 1$, and consequently the exact
first character of the tail difference is the $\times 2$ orbit of one real
number:
\[
  \mathrm{tailOrbitFirstExp}(h,N)
  \;=\; e\bigl(2^N\alpha_h\bigr),
  \qquad e(x):=\exp(2\pi i x).
\]
Hence, for fixed $h$, the phases $\{\,R_{N+h}-R_N \bmod 1\,\}_{N\ge 0}$ are
the forward orbit of $\alpha_h \bmod 1$ under $x\mapsto 2x$.
\end{lem}

\noindent
The second identity and the character form are landed:
\lean{Erdos249257.\allowbreak TotientTailPeriodKiller.\allowbreak tail_diff_eq_scaled_\allowbreak totient_series_sub_prefix}{PivotAntiReconstruction.lean:54}
and
\lean{Erdos249257.\allowbreak TotientTailPeriodKiller.\allowbreak tailOrbitFirstExp_eq_scaledTotientSeriesFirstExp}{PivotAntiReconstruction.lean:66}
\quad \ev{Lean} \quad \scale{uniform} \quad \coord{other:binary-window}.
The recurrence $R_{N+1}=2R_N-\varphi(N+1)$ is one line from the definition
and is used below only for intuition \ev{Math}.

This is not a new fact --- it is two landed Lean theorems read together ---
but reading them together has a consequence the corpus never draws. The
first-harmonic block sum, which is the analytic frontier of \#249, is a
\emph{lacunary exponential sum in a single real variable}. The following
proposition makes the transfer exact, including the truncation budget.

\begin{prop}[Orbit transfer]\label{prop:transfer}
Define
\[
  \mathrm{OrbitBlockGap} :\equiv
  \forall h\ge 1\ \forall X_0\ \exists X\ge \max(X_0,1):\quad
  \sum_{N=X}^{2X-1}\cos\bigl(2\pi\,2^{N}\alpha_h\bigr)\;\le\;\tfrac{89}{100}\,X .
\]
Then $\mathrm{OrbitBlockGap}$ implies that $S$ is irrational.
\end{prop}

\begin{proof}
Fix $h\ge 1$ and $X_0$, and take $X$ as supplied. Choose $L$ with
$2^{L}\ge 1024\,(2X+h+L+2)$; such $L$ exists because $2^L/L\to\infty$, and it
satisfies the room inequality $16(2X+h+L+2)\le 2^L$ a fortiori. For every
$N<2X$ the landed truncation bound
\lean{Erdos249257.\allowbreak TotientTailPeriodKiller.\allowbreak norm_windowFirstExp_\allowbreak sub_tailOrbitFirstExp_lt}{PivotAntiReconstruction.lean:202}
gives
$\bigl\|\,\mathrm{windowFirstExp}(h,N,L)-e(2^N\alpha_h)\,\bigr\|
 < 2\pi(N+L+h+2)/2^{L} \le 2\pi/1024 < 1/100$,
so by
\lean{Erdos249257.\allowbreak TotientTailPeriodKiller.\allowbreak windowFirstCos_le_\allowbreak add_of_tailOrbitGap}{PivotAntiReconstruction.lean:220},
$\mathrm{windowFirstCos}(h,N,L)\le \cos(2\pi 2^N\alpha_h)+1/100$. Summing over
the $X$ values $N\in[X,2X)$ gives
$\sum_N \mathrm{windowFirstCos}(h,N,L)\le \tfrac{89}{100}X+\tfrac{1}{100}X
=\tfrac{9}{10}X$, so
\lean{Erdos249257.\allowbreak TotientTailPeriodKiller.\allowbreak exists_certifiedKill_of_\allowbreak first_harmonic_gap}{FirstHarmonicGap.lean:128}
produces $N\in[X,2X)$ with $\mathrm{certifiedKill}(h,N,L)$, and $N\ge X\ge X_0$.
This is exactly the certificate supply consumed by
\lean{Erdos249257.\allowbreak TotientTailPeriodKiller.\allowbreak irrational_totient_series_\allowbreak of_certificate_supply}{TotientTailPeriodKiller.lean:394}.
\end{proof}

\noindent
\ev{Math} \quad \scale{cofinal} \quad \coord{other:binary-window}. The proof
uses only landed lemmas; formalising it is a variant of
\lean{Erdos249257.\allowbreak irrational_totient_series_of_\allowbreak first_harmonic_norm_gap}{FirstHarmonicPivot.lean:83}
with the constant $21/25$ replaced by the real-part constant $89/100$ and the
room factor $16$ by $1024$. Nothing here is an improvement of the analytic
requirement; it is a change of coordinate that makes the requirement legible.

\begin{cor}[Digit form of the analytic requirement]\label{cor:digitform}
Let $\rho_h(X)$ denote the proportion of $N\in[X,2X)$ with
$\|2^N\alpha_h\|_{\R/\Z}\ge 1/4$. If for every $h\ge 1$ there are cofinally
many $X$ with $\rho_h(X)\ge 11/100$, then $S$ is irrational.
Equivalently, when $\alpha_h$ is not a dyadic rational: if for every $h$
there are cofinally many dyadic blocks $[X,2X)$ in which the binary
expansion of $\alpha_h$ contains at least $0.11X$ digit changes --- that is,
in which the mean binary run length is at most $9.1$ --- then $S$ is
irrational.
\end{cor}

\begin{proof}
$\|x\|\ge 1/4$ forces $\cos 2\pi x\le 0$, and the remaining terms are at most
$1$, so $\sum_{N}\cos(2\pi 2^N\alpha_h)\le (1-\rho_h(X))X$; require
$1-\rho\le 89/100$ and apply Proposition~\ref{prop:transfer}. For the digit
reading, $\|2^N\alpha_h\|\ge 1/4$ holds exactly when the binary digits of
$\alpha_h$ in positions $N+1,N+2$ differ.
\end{proof}

\noindent
\ev{Math}. This threshold is worth staring at. A single digit change in
every ninth position, in one block per scale, for each fixed $h$, would
settle Erd\H{o}s \#249. What is actually true, numerically, is far stronger.

\begin{table}[h]\centering
\begin{tabular}{r r r r r}
\hline
$h$ & $X$ & $\rho_h(X)$ & $X^{-1}\sum_{N\in[X,2X)}\cos(2\pi 2^N\alpha_h)$ & required \\
\hline
$1$ & $256$  & $0.4766$ & $\phantom{-}0.0037$ & $\le 0.89$ \\
$1$ & $2048$ & $0.5073$ & $-0.0043$ & $\le 0.89$ \\
$3$ & $256$  & $0.4648$ & $\phantom{-}0.0296$ & $\le 0.89$ \\
$3$ & $2048$ & $0.5020$ & $\phantom{-}0.0027$ & $\le 0.89$ \\
$8$ & $2048$ & $0.5112$ & $-0.0088$ & $\le 0.89$ \\
$13$& $2048$ & $0.5049$ & $-0.0195$ & $\le 0.89$ \\
\hline
\end{tabular}
\caption{Block statistics for $\alpha_h=(2^h-1)S$, computed from the exact
integer $\lfloor 2^{6000}S\rfloor$ (error $<1$ in the last bit, tail bound
$(M+2)2^{-M}$ at $M=6080$). Over positions $1..5800$ the binary expansion of
$\alpha_1$ has $2864$ runs, mean run length $2.0251$, longest run $13$; for
$\alpha_3$, $2964$ runs, mean $1.9568$, longest $12$; for $\alpha_8$, $2927$
runs, mean $1.9816$, longest $14$. A uniform digit model predicts mean run
length $2$ and longest run $\approx\log_2 5800\approx 12.5$. \ev{Cert}}
\label{tab:blocks}
\end{table}

\begin{obs}\label{obs:allroutes}
Under Lemma~\ref{lem:orbit} every route that survives the barrier
classification becomes a statement about the binary expansion of $S$ or of
some $\alpha_h$. Routes 1 and 2 ask for a positive proportion of digit
changes in a block; Route 3 asks for one short run at a prime-indexed
position; Route 4 asks that two blocks of digits of $S$ at distance $h$
disagree early. The routes differ in \emph{which positions} they interrogate
and \emph{how uniform} the margin must be, not in what they are about. This
is developed in \S\ref{sub:shape}.
\end{obs}

\subsection{Route 1: a constant-saving cancellation over a dyadic block}

\paragraph{The exact statement needed.}
\[
  \forall h\ge 1\ \forall X_0\ \exists X,L:\quad
  \max(X_0,1)\le X,\quad 16(2X+h+L+2)\le 2^{L},\quad
  \Bigl\|\sum_{N=X}^{2X-1} e\bigl(D(h,N,L)/2^{L}\bigr)\Bigr\|\le \tfrac{21}{25}X,
\]
where $D(h,N,L)=\sum_{j<L}\bigl(\varphi(N+h+1+j)-\varphi(N+1+j)\bigr)2^{L-1-j}$.
\lean{Erdos249257.\allowbreak TotientTailPeriodKiller.\allowbreak DTWFirstHarmonicNormGap}{FirstHarmonicPivot.lean:75},
consumer
\lean{Erdos249257.\allowbreak irrational_totient_series_of_\allowbreak first_harmonic_norm_gap}{FirstHarmonicPivot.lean:83}
\quad \ev{Open} \quad \scale{cofinal}. By Proposition~\ref{prop:transfer} the
weaker real-part form with constant $89/100$ suffices.

\paragraph{What $D$ is, as an arithmetic object.}
Dividing by $2^L$ and reindexing the two ranges onto a common variable
$t$, with $m=N+t$,
\begin{equation}\label{eq:window}
  \frac{D(h,N,L)}{2^{L}}
  =\underbrace{-\sum_{t=1}^{h}\frac{\varphi(N+t)}{2^{t}}}_{\text{head}}
  \;+\;\underbrace{(2^{h}-1)\sum_{t=h+1}^{L}\frac{\varphi(N+t)}{2^{t}}}_{\text{body}}
  \;+\;\underbrace{\sum_{t=L+1}^{L+h}\frac{\varphi(N+t)}{2^{t-h}}}_{\text{tail}} .
\end{equation}
So the phase is a \emph{geometrically weighted linear form in $\varphi$ at
$L+h$ consecutive arguments}, with the value at offset $t$ entering modulo
$2^{t}$ and with an odd multiplier $2^h-1$ throughout the body. Two features
control everything that follows. First, the weight at offset $t$ has
denominator exactly $2^{t}$, so offset $t$ contributes nothing unless
$v_2(\varphi(N+t))<t$: the $2$-adic valuation of the totient is the gate.
Second, $2^h-1$ is odd, so the multiplier never closes a gate that
$\varphi$ leaves open.

\paragraph{Which technique family is even the right shape.}
Four candidates, in decreasing order of relevance.

\emph{Weyl differencing and van der Corput} are the wrong shape and can be
dismissed exactly. Both require the phase to be a smooth or polynomial
function of the summation variable so that a difference operator lowers its
degree. Here $N\mapsto D(h,N,L)/2^{L}$ is, by Lemma~\ref{lem:orbit}, the
$\times 2$ orbit map: differencing in $N$ multiplies the phase by $2$ and
subtracts an integer, so the difference operator is an exact isometry of the
problem and lowers nothing. This is not a heuristic --- it is
$R_{N+1}=2R_N-\varphi(N+1)$ read modulo $1$.

\emph{The large sieve} needs a family of well-separated frequencies and a
sum over both the frequencies and the variable. Here there is one frequency
per $N$ and no family: the sum is over a single orbit. A large-sieve
inequality could be applied after introducing an artificial family (for
example, over the $h$ parameter), but the predicate quantifies $h$
universally, so an average over $h$ is not admissible.

\emph{Vaughan/Vinogradov bilinear decomposition} is the right shape, and is
what Route 2 already implements: decompose the argument $N+t$ at a large
prime factor, use $\varphi(mp)=\varphi(m)(p-1)$ to linearise the phase in
$p$, and sum over primes. \S\ref{sub:pivot} carries this as far as it goes.

\emph{Erd\H{o}s's 1948 digit method}, which proved irrationality of
$E=\sum_n 1/(2^n-1)=\sum_N d(N)/2^N$, is the right shape for the
\coord{mobius-mersenne} coordinate. That coordinate exists here and is
exact: from $\varphi=\mu * \mathrm{id}$,
\begin{equation}\label{eq:mobmers}
  S=\sum_{d\ge 1}\mu(d)\,\frac{2^{d}}{(2^{d}-1)^{2}}
  \qquad (\ev{Math};\ \text{verified to }16\ \text{digits},\ S=1.3676308019850223\ldots,\ \ev{Cert}).
\end{equation}
But the method does not transfer, for a reason that can be stated
quantitatively rather than vaguely. Erd\H{o}s's argument for $E$ exploits
that the coefficient $d(N)$ has enormous multiplicative fluctuation: over
$N\le Y$ the ratio of its maximum to its typical value is
$Y^{(\log 2+o(1))/\log\log Y}$, so certain highly composite $N$ deposit
identifiable spikes into the digit stream. The coefficient here is
$\varphi(N)$, which satisfies $\varphi(N)/N\in[c/\log\log N,1]$: its
multiplicative fluctuation over $N\le Y$ is a factor $e^{\gamma}\log\log Y$,
that is, $\log\log Y$ rather than $Y^{c/\log\log Y}$. There are no spikes to
find. In the M\"obius--Mersenne coordinate \eqref{eq:mobmers} the situation is
worse rather than better: the coefficients $\mu(d)$ change sign, so the
digit blocks contributed by successive $d$ cancel rather than accumulate,
and there is no positivity to run a block argument on. This is, as far as
the evidence in this corpus goes, the specific reason \#249 is harder than
its Erd\H{o}s--Borwein ancestor.

\paragraph{The specific obstacle: what the estimate cannot be deduced from.}
The corpus records that no theorem proves the first-harmonic gap
inequivalent to irrationality, and argues the point from the structure of a
collapse proof. It can be settled outright, and the witness lives in the
same coefficient class that carries barriers B1 and B7.

\begin{thm}[The gap is strictly stronger than irrationality]\label{thm:lacunary}
Let $c(n)=1$ if $n=k!$ for some $k\ge 1$ and $c(n)=0$ otherwise, so
$0\le c(n)\le n$ for all $n\ge 1$, and let $\beta=\sum_{n\ge1}c(n)/2^{n}
=\sum_{k\ge 1}2^{-k!}$. Then $\beta$ is irrational, and for every $h\ge 1$
and every $X\ge 81(h+5)$,
\[
  \sum_{N=X}^{2X-1}\cos\bigl(2\pi\,2^{N}(2^{h}-1)\beta\bigr) \;>\; \tfrac{9}{10}X .
\]
Consequently the block-gap requirement fails at \emph{every} scale for
$\beta$, while $\beta$ is irrational. No proof of the block gap for $S$ can
therefore proceed from the irrationality of $S$ together with the growth
bound $c(n)\le n$; it must use arithmetic of $\varphi$.
\end{thm}

\begin{proof}
Irrationality of $\beta$ is Liouville's criterion. Write
$\gamma=(2^{h}-1)\beta=\sum_{k}(2^{h-k!}-2^{-k!})$, so for every $k$ with
$k!>2h$ the binary expansion of $\gamma$ carries ones exactly in the
positions $k!-h+1,\dots,k!$ and zeros between consecutive such blocks. If
$N\le k!-h-5$ for the least $k$ with $k!>N$, then
$\{2^{N}\gamma\}\le 2\cdot 2^{-5}=1/16$, whence $\|2^{N}\gamma\|\le 1/16$ and
$\cos(2\pi 2^{N}\gamma)\ge\cos(\pi/8)>0.9239$. The remaining $N$ lie in
$(k!-h-5,\,k!]$, at most $h+5$ values per factorial, and consecutive
factorials differ by a factor exceeding $2$ from $k\ge 3$, so the window
$[X,2X)$ meets at most one such interval. Therefore the sum exceeds
$0.9239\,(X-h-5)-(h+5)$, which exceeds $\tfrac{9}{10}X$ once
$X\ge 81(h+5)$.
\end{proof}

\noindent
\ev{Math} \quad \scale{cofinal} \quad \coord{other:binary-window}. This
upgrades the corpus's ``no collapse mechanism is known to apply'' to ``no
collapse mechanism can exist''. It also says precisely what the missing
input must do: it must rule out that $S$ behaves, in its doubling orbit,
like a Liouville number. That is the content of Corollary~\ref{cor:digitform},
and it is why the requirement is a digit-statistics statement and not an
irrationality statement.

\paragraph{Size of the quantity, honestly calibrated.}
Under any model in which the binary digits of $\alpha_h$ behave like fair
coin flips, $\rho_h(X)\to 1/2$ and the block sum is $O(\sqrt{X\log\log X})$
by the Erd\H{o}s--G\'al law of the iterated logarithm for lacunary series
\ev{Cited}. The requirement is $\rho_h\ge 0.11$ and a saving of a constant
factor. The measured values in Table~\ref{tab:blocks} are $\rho_h\approx 0.50$ and a block
average of order $X^{-1/2}$. So the analytic requirement is weaker than the
apparent truth by a factor $\sqrt X$, and the difficulty is entirely that no
technique produces \emph{any} nontrivial digit statistic for an explicit
constant of this kind. That is a statement about technology, and it is worth
separating from a statement about $S$.

\subsection{Route 2: the four-term pivot budget}\label{sub:pivot}

\paragraph{The exact statement needed.}
For every $h\ge 1$ there must exist $s\ge 1$ and $\eta\in(0,1)$ such that for
every $X_0$ there are $X\ge\max(X_0,1)$ and $L$ with $h\le L-s$,
$16(2X+h+L+2)\le 2^{L}$, and all four of
\[
\begin{aligned}
  \operatorname{Re}\bigl(\mathrm{pivotCenteredCorrelation}\bigr) &\le \tfrac{14}{25}X,
  &\qquad \bigl\|\mathrm{pivotFiberMeanContribution}\bigr\| &\le \tfrac{1}{100}X,\\
  \bigl\|\mathrm{pivotBadContribution}\bigr\| &\le \tfrac{1}{100}X,
  &\qquad \bigl\|\mathrm{pivotNonSupplierContribution}\bigr\| &\le \tfrac{8}{25}X .
\end{aligned}
\]
\lean{Erdos249257.\allowbreak TotientTailPeriodKiller.\allowbreak DTWPivotResidualDecorrelation}{FirstHarmonicPivot.lean:397},
consumer
\lean{Erdos249257.\allowbreak irrational_totient_\allowbreak series_of_pivotResidualDecorrelation}{FirstHarmonicPivot.lean:405}
\quad \ev{Open} \quad \scale{cofinal}. The decomposition is an exact identity,
\lean{Erdos249257.\allowbreak TotientTailPeriodKiller.\allowbreak windowFirstExp_sum_\allowbreak eq_pivot_decomposition}{FirstHarmonicPivot.lean:342}
\quad \ev{Lean}, and the supplier set at the canonical fibre is a membership
equality with a shifted dyadic interval of primes,
\lean{Erdos249257.\allowbreak TotientTailPeriodKiller.\allowbreak mem_pivotFiber_\allowbreak one_overlap_iff}{FirstHarmonicPivot.lean:251}
\quad \ev{Lean}.

\paragraph{One budget has a classical one-sided majorant.}
The non-supplier term is not identified exactly by smooth numbers.  What the
factorisation gives is the one-way containment needed for an upper bound.

\begin{prop}[Non-suppliers: the valid one-sided estimate]\label{prop:dickman}
Choose the admissible depth $L$ minimally for each large $X$, with the
predicate's previously fixed $h$ and $s$.  If
$n=N+L-s+1$ is not a supplier, then its largest prime factor satisfies
\[
  P(n)\le (4+o(1))\sqrt X .
\]
Therefore
\[
  \#\{N\in[X,2X):\neg\,\mathrm{pivotSupplier}(X,L,s,N)\}
  \le
  \Psi(2X+O(\log X),(4+o(1))\sqrt X)
  -\Psi(X+O(\log X),(4+o(1))\sqrt X).
\]
A uniform smooth-number asymptotic for this shifted interval would make the
right side $(1-\log2+o(1))X<\tfrac8{25}X$ and would discharge the budget.
\end{prop}

\begin{proof}
Let $p=P(n)$ and $m=n/p$.  Supplier failure means either
$p\le2\lfloor\sqrt X\rfloor$ or $m>\lfloor\sqrt X\rfloor/2$ (up to the
definition's exact floor and positivity clauses).  In the second case
$p=n/m\le(4+o(1))\sqrt X$, because minimal admissible $L$ gives
$L-s+1=O(\log X)$ and hence $n<2X+O(\log X)$.  Thus every non-supplier lies in
the displayed smooth-number set.  This proves only an upper bound, not an
asymptotic equality for the non-supplier count.
\end{proof}

\noindent
\ev{Math} for the containment; \ev{Cited} for the smooth-number estimate
needed to finish the numerical budget.  In particular the earlier identity
with Dickman density, and numerical extrapolations of an exact crossover
scale, are not claimed.  The implication is sufficient if the required
uniform shifted-interval estimate is supplied.

\begin{prop}[The bad-cofactor budget]\label{prop:badcof}
For $\eta\in(0,1)$ let $B(\eta)=\{m:\varphi(m)<\eta m\}$, with natural density
$D(\eta)$; by Schoenberg's theorem $D$ exists, is continuous, and
$D(0+)=0$ \ev{Cited}. Then
\[
  \#\{N\in[X,2X): N\ \text{a supplier with cofactor}\ m\in B(\eta)\}
  \;\le\;\bigl(D(\eta)+o(1)\bigr)X ,
\]
so a single choice of $\eta$ with $D(\eta)<1/200$ meets the $\tfrac{1}{100}X$
budget for all large $X$. Since $\eta$ is quantified existentially before
$X_0$, and a smaller $\eta$ only enlarges the good set, this choice costs
nothing elsewhere in the budget.
\end{prop}

\begin{proof}[Proof sketch]
Suppliers with cofactor $m$ number $\pi(2X/m)-\pi(X/m)$, which is at most
$(2X/m)/\log X$ up to $1+o(1)$ because $m\le\sqrt X/2$ forces
$\log(X/m)\ge\tfrac12\log X$. Sum over $m\in B(\eta)$ with $m\le\sqrt X/2$
and use that the logarithmic density of $B(\eta)$ tends to $D(\eta)$.
\end{proof}

\noindent
\ev{Math} \quad \scale{cofinal}. The choice is not tight: the least element
of $B(1/5)$ is $2\cdot3\cdot5\cdot7\cdot11\cdot13=30030$, with
$\varphi(m)/m=5760/30030=0.19181$, and the logarithmic density of $B(1/5)$ up
to $3\cdot10^{5}$ is $3\cdot 10^{-5}$ \ev{Cert}. So $\eta=1/5$ is already
defensible pending an explicit numerical bound on $D(1/5)$, and that bound is
itself a finite computation a reader could start today.

\paragraph{The pivot quantifiers force a deep modulus.}
The pivot sits at offset $t=L-s+1$ in \eqref{eq:window}.  Crucially, the
definition of \texttt{DTWPivotResidualDecorrelation} chooses $s$ once, after
$h$ and before the universal threshold $X_0$.  Thus $s$ is fixed as
$X\to\infty$; it cannot be increased with $X$ to make $t$ shallow.  The room
condition gives
\[
  2^L\ge16(2X+h+L+2),\qquad
  t=L-s+1\ge\log_2X-O_{h,s}(1).
\]
After removing the $2$-part of the coefficient, the prime-progression modulus
is still typically of order $X$ up to subpolynomial factors, not
$(\log X)^A$.  Standard Siegel--Walfisz therefore does not reach the
quantifier order of the stated predicate; choosing a hypothetical
$t=O(\log\log X)$ would amount to choosing $s$ after $X$, which the predicate
forbids.

The exact $2$-adic triviality observation remains useful: if
$v_2(\varphi(m))\ge t$, then the pivot phase is $1$.  It does not create a
nonempty analytic window under the actual quantifiers.  Hence the claimed
shallow-modulus route is withdrawn.  This does not disprove the pivot
predicate; it says only that the proposed prime-distribution argument does
not supply it. \ev{Math} \scale{cofinal}.

\paragraph{The specific obstacle: the weight is neither Type I nor Type II.}
Propositions~\ref{prop:dickman} and \ref{prop:badcof} isolate possible
bookkeeping bounds, subject to their stated uniform estimates.  The
fixed-$s$ quantifier leaves the fibre-mean modulus deep, and the centred term
has a separate correlation obstruction.  Factoring the pivot argument as
$n=mp$ and dividing out the pivot phase leaves
\[
  \operatorname{Re}\sum_{m}\ \sum_{p\,:\,mp\in[X,2X)+t}
    w(mp)\,\bigl(e\bigl(a_m(p-1)/2^{t}\bigr)-\overline{e}_m\bigr)
  \;\le\;\tfrac{14}{25}X ,
  \qquad |w|\equiv 1 ,
\]
where $\overline{e}_m$ is the fibre mean and, by \eqref{eq:window}, $w(n)$ is
the character of the same weighted totient sum with the single offset $t$
deleted. Vinogradov's bilinear method requires that, after the decomposition
$n=mp$, the summand split as $\alpha_m\beta_p$ (Type II) or be independent of
$p$ given $m$ (Type I), up to boundedly many pieces. The weight $w(mp)$ is
neither: it is a function of the totients of the $L+h$ integers
\emph{adjacent} to $mp$, and $\varphi(mp+j)$ for $j\ne0$ bears no relation to
the factorisation $mp$ --- it is not a function of $m$, not a function of
$p$, and not a product of the two. No device is known that converts a weight
of this shape into bilinear form.

At $h=1$ and the canonical pivot the required estimate becomes fully
explicit. There $e\bigl((p-1)/4\bigr)=i^{\,p-1}=\chi_{-4}(p)$, so what is
needed is
\[
  \operatorname{Re}\sum_{X<p\le 2X}\chi_{-4}(p)\,
    e\Bigl(\frac{\varphi(p+1)}{8}+\frac{\varphi(p+2)}{16}+\cdots\Bigr)
  \;\le\;\Bigl(\frac{9}{10}-\delta\Bigr)\,\bigl(\pi(2X)-\pi(X)\bigr).
\]
The missing input is thus a \emph{non-correlation between a fixed quadratic
character at $p$ and the $2$-adic behaviour of $\varphi$ at the shifts
$p+1,p+2,\dots$}: a correlation statement for two arithmetic functions at
shifted arguments, of Chowla--Elliott type. The unconditional results in that
family --- Matom\"aki--Radziwi\l{}\l{} in almost all short intervals, Tao's
logarithmically averaged Chowla, Tao--Ter\"av\"ainen for odd order
\ev{Cited} --- are averaged, never pointwise at a single scale.

\paragraph{A swing: the predicate needs only cofinally many $X$, so average
over $X$.} This is the one place where the shape of the obligation is a gift.
Both $\mathrm{DTWFirstHarmonicNormGap}$ and
$\mathrm{DTWPivotResidualDecorrelation}$ ask for \emph{some} $X$ beyond each
$X_0$, never for all $X$. Hence it suffices to prove a logarithmically
averaged bound: if
\[
  \sum_{k\le K}\ \frac{1}{2^{k}}
  \Bigl|\sum_{N\in[2^{k},2^{k+1})}e\bigl(D(h,N,L_k)/2^{L_k}\bigr)\Bigr|
  \;=\;o(K)
  \quad\text{for each }h,
\]
with $L_k$ any admissible depth sequence, then infinitely many blocks satisfy
the gap and Proposition~\ref{prop:transfer} applies. Logarithmic averaging is
exactly the regime in which the entropy-decrement method operates. The first
concrete step is not to prove this but to decide whether the method reaches
it: the phase $e\bigl(D(h,N,L)/2^{L}\bigr)$ is not multiplicative, but it is
a bounded \emph{local function of a multiplicative function} --- a fixed
continuous function of $\bigl(\varphi(N+1)/2,\dots,\varphi(N+L)/2^{L}\bigr)$
read modulo $1$ --- whereas entropy decrement is stated for correlations of
bounded multiplicative functions along fixed shifts. The precise question to
settle, and it can be started on immediately, is whether the decrement
survives when the observable is a local function of $\varphi$ with a growing
number of coordinates $L\approx\log_2X$ rather than a product of boundedly
many multiplicative values. \ev{Open}

\subsection{Route 3: uniform quantitative escape at primes}

\paragraph{The exact statement needed.}
\[
  \forall h\ge 1\ \forall N_0\ \exists p\ \text{prime}:\quad
  \max(N_0+h+1,\,h+5)\le p
  \ \wedge\
  \operatorname{Re}\bigl(\mathrm{tailOrbitFirstExp}(h,\,p-h-1)\bigr)<\tfrac{9}{10}.
\]
\lean{ErdosProblems.\allowbreak Erdos249.\allowbreak DTWNaturalPrimeTailOrbitStrictGap}{ErdosProblems/Erdos249/TotientStrictPrimeEscape.lean:157},
consumer
\lean{ErdosProblems.\allowbreak Erdos249.\allowbreak irrational_totient_\allowbreak series_of_naturalPrimeTailOrbitStrictGap}{ErdosProblems/Erdos249/TotientStrictPrimeEscape.lean:205}
\quad \ev{Open} \quad \scale{cofinal}.

\paragraph{What it says, after Lemma~\ref{lem:orbit}.}
$\operatorname{Re}e(2^{p-h-1}\alpha_h)<9/10$ says
$\|2^{p-h-1}\alpha_h\|_{\R/\Z}>\arccos(9/10)/2\pi=0.0717831\ldots$, so the
requirement is exactly: \emph{for each $h$, the indicated prime-indexed
$\times2$ orbit point stays more than $0.0717831\ldots$ from the nearest
integer, for cofinally many primes $p$}.  A short-run condition on the binary
digits can be a sufficient proxy for this circle-distance inequality, but it
is not equivalent to it; a run-length description alone does not determine
the residual position inside the dyadic cylinder.

\paragraph{The specific obstacle.} Two, stacked. First, this is a
\emph{pointwise} demand: it names an index, and the promotion audit records
that no pointwise producer in this programme has ever been supplied at even
one large index, across seven independent attempts, while the single audited
row whose missing input is a block average is the first-harmonic row.
Second, and sharper: digit behaviour \emph{along the primes} is strictly
harder than digit behaviour on average even for constants where the average
case is settled. The only constants whose digits are understood at all along a sparse
subsequence are ones built for the purpose (Champernowne;
Copeland--Erd\H{o}s) \ev{Cited}. No technique reads the digits of an
explicit constant at prime positions.

\paragraph{Numerical calibration.} Among the $501$ primes in $[1000,5000)$,
the proportion with $\operatorname{Re}e(2^{p-h-1}\alpha_h)<9/10$ is $0.8623$
for $h=1$, $0.8503$ for $h=2$, $0.8583$ for $h=5$, with minima below
$-0.9999$ \ev{Cert}. The cofinal supply is thus met by about six primes in
seven at these scales. As with Route 1, the requirement is very weak and the
obstacle is that nothing sees it.

\paragraph{A swing: replace the pointwise demand by a canonical fibre.}
The corpus warns against the subset consumer
\lean{Erdos249257.\allowbreak TotientTailPeriodKiller.\allowbreak exists_certifiedKill_of_\allowbreak first_harmonic_gap_subset}{FirstHarmonicGap.lean:104}
because a predicate quantifying existentially over subsets
$T\subseteq[X,2X)$ is collapse-exposed: the proved collapse
\lean{Erdos249257.\allowbreak dtwWindowSeparatedPairs_iff_\allowbreak irrational_totient_series}{PivotAntiReconstruction.lean:1765}
works precisely by \emph{choosing} $T=\{N,N+1\}$. That warning applies to a
free $T$. It does not apply to a $T$ named in advance as a function of
$(h,X,L)$, and the corpus already contains the right one:
\[
  T_{h,X,L}\;:=\;\mathrm{pivotFiber}(X,L,L-h,1)
  \;=\;\{\,N\in[X,2X)\ :\ N+h+1\ \text{prime}\,\},
\]
a membership equality rather than a sampled surrogate
(\lean{Erdos249257.\allowbreak TotientTailPeriodKiller.\allowbreak mem_pivotFiber_\allowbreak one_overlap_iff}{FirstHarmonicPivot.lean:251},
\ev{Lean}). The proposal is to prove
\begin{equation}\label{eq:primefibre}
  \forall h\ge1\ \forall X_0\ \exists X\ge X_0,\ L:\quad
  16(2X+h+L+2)\le 2^{L}
  \ \wedge\
  \sum_{N\in T_{h,X,L}}\mathrm{windowFirstCos}(h,N,L)\ \le\ \tfrac{9}{10}\,\bigl|T_{h,X,L}\bigr| ,
\end{equation}
which by the landed subset consumer yields a certificate at some
$N\ge X\ge X_0$, hence irrationality. \ev{Open}

Three features make this the sharpest available target. (i) On this fibre the
cofactor is $m=1$, so the pivot coefficient is $(2^{h}-1)\varphi(1)=2^{h}-1$,
which is \emph{odd}, so the $2$-adic triviality described above cannot occur.
(ii) The pivot modulus is
$2^{h+1}$, \emph{fixed} once $h$ is fixed, so the fibre-mean term is governed
by the prime number theorem in progressions to a fixed power of two, where
the Ramanujan main term $c_{2^{h+1}}(2^{h}-1)=\mu(2^{h+1})=0$ vanishes for
every $h\ge1$ and the error term is effective; the fixed-$s$ deep-modulus
obstruction does not apply to this separately defined canonical fibre.
(iii) The bookkeeping budgets of \S\ref{sub:pivot} disappear:
inside $T$ there are no non-suppliers and no bad cofactors. What remains is
exactly one estimate --- decorrelation of the fixed phase
$e\bigl((2^{h}-1)(p-1)/2^{h+1}\bigr)$ from the residual totient weight --- and
it is the same estimate isolated at the end of \S\ref{sub:pivot}. This does
not make that estimate easier. It removes everything that is not that
estimate.

\begin{rem}
Honesty about \eqref{eq:primefibre}: it is not proved to be strictly stronger
than irrationality, and no collapse is proved either. What can be said is
that the one collapse mechanism that exists in this lane consumes the freedom
to choose the sample, and \eqref{eq:primefibre} has no such freedom. Deciding
this either way --- exhibiting a collapse, or a witness in the coefficient
class of Theorem~\ref{thm:lacunary} that separates it from irrationality ---
is itself a well-defined finite piece of work.
\end{rem}

\subsection{Route 4: depth-locked full-depth escape}

\paragraph{The exact statement needed.}
\[
  \mathrm{ApFullDepthEscape} :\equiv
  \forall d\ge 1\ \forall N\ \exists t\ge 1:\ \mathrm{certifiedKill}(td,\,N,\,td),
\]
unpacked, $(N+2td+2:\Z)<D(td,N,td)\bmod 2^{td}<2^{td}-(N+2td+2)$.
\lean{ErdosProblems.\allowbreak Erdos249.\allowbreak PeriodMultipleEscape.\allowbreak ApFullDepthEscape}{ErdosProblems/Erdos249/PeriodMultipleEscape.lean:162},
consumer
\lean{ErdosProblems.\allowbreak Erdos249.\allowbreak PeriodMultipleEscape.\allowbreak irrational_totient_\allowbreak series_of_apFullDepthEscape}{ErdosProblems/Erdos249/PeriodMultipleEscape.lean:171},
ambient equivalence
\lean{ErdosProblems.\allowbreak Erdos249.\allowbreak PeriodMultipleEscape.\allowbreak periodMultipleKillSupply_iff_irrational}{ErdosProblems/Erdos249/PeriodMultipleEscape.lean:153}
\quad \ev{Open} \quad \scale{cofinal}. It is the shortest fully stated open
inequality the programme has produced.

\paragraph{What it says: an anti-self-similarity statement about the digits
of $S$ alone.}
Writing $h=td$ and using
$D(h,N,L)/2^{L}=(R_{N+h}-R_N)-(R_{N+L+h}-R_{N+L})/2^{L}$ at $L=h$, together
with the landed strip $|R_{M+h}-R_M|<M+h+2$
(\lean{Erdos249257.\allowbreak TotientTailPeriodKiller.\allowbreak abs_tail_\allowbreak diff_lt}{CarrySurvivorExtinction.lean:79},
\ev{Lean}), applied at $M=N+h$, one gets:

\begin{prop}\label{prop:route4}
$\mathrm{certifiedKill}(h,N,h)$ holds whenever
$\bigl\|2^{N+h}S-2^{N}S\bigr\|_{\R/\Z}>2(N+2h+2)/2^{h}$.
\end{prop}

\begin{proof}
$\|D(h,N,h)/2^{h}\|\ge\|R_{N+h}-R_N\|-|R_{N+2h}-R_{N+h}|/2^{h}
 >2(N+2h+2)/2^{h}-(N+2h+2)/2^{h}=(N+2h+2)/2^{h}$, and
$R_{N+h}-R_N\equiv 2^{N}(2^{h}-1)S=2^{N+h}S-2^{N}S\pmod 1$ by
Lemma~\ref{lem:orbit}. A residue at distance more than $(N+2h+2)/2^{h}$ from
$\Z$ is exactly a certified kill at depth $h$.
\end{proof}

\noindent
\ev{Math}. So Route 4 asks: \emph{for every ray $d$ and every basepoint $N$,
some multiple $h=td$ is such that the tail of the binary expansion of $S$
beginning at position $N+1$ and the tail beginning at position $N+h+1$, read
as reals in $[0,1)$, are more than $2(N+2h+2)/2^{h}$ apart in $\R/\Z$.} Two
things follow. First, unlike Routes 1--3, this involves only $S$ itself, not
the multiples $\alpha_h$. Second, the required precision grows only like
$\log_2(N+2h)$ while the available depth grows linearly in $t$: the demand
weakens rapidly along each ray.

\paragraph{Which technique family fits, and why it is unavailable.}
A statement that a sequence does not almost repeat at some multiple of every
gap, to logarithmic precision, is a \emph{subword-complexity} statement, and
the mature technology for such statements about a specific real number is the
Adamczewski--Bugeaud method: the Schmidt subspace theorem applied to the
$b$-ary expansion, which shows that an algebraic irrational cannot have too
many long repetitions \ev{Cited}. The shape fits exactly. The hypothesis does
not: the method conditions on algebraicity, and $S$ is not known to be
algebraic --- indeed it is not known to be irrational, which is the whole
problem. There is no version of the method that runs from an analytic
description of the constant. That is the entire obstruction, and it is worth
stating plainly because it explains why Route 4 looks so much easier than it
is: the required combinatorial property is weak, but the only machine that
produces such properties needs an input we do not have.

\paragraph{What the data says.}
At $N=300$, the least $t$ with $\mathrm{certifiedKill}(td,300,td)$ exists
and is small for every ray $d\le 24$: $t=10$ for $d=1$, $t=5$ for $d=2$,
$t=4$ for $d=3$, $t=3$ for $d=4$, $t=2$ for $5\le d\le 9$, and $t=1$ for
$10\le d\le 24$ --- in every case at or near the first depth admitted by the
room floor. The residues are not marginal: $r/2^{h}=0.3594$ at $d=1$,
$0.5977$ at $d=3$, $0.6853$ at $d=7$, $0.2620$ at $d=17$, $0.5887$ at
$d=23$, against edge radii $3.1\cdot10^{-1}$, $8.0\cdot10^{-2}$,
$2.0\cdot10^{-2}$, $2.6\cdot10^{-3}$, $4.1\cdot10^{-5}$ \ev{Cert}. This is a
floor, not a trend, exactly as barrier B1 requires; it decides nothing.

\paragraph{A heuristic size computation.}
Under a uniform model for $\|2^{N}(2^{h}-1)S\|$ the failure probability at
depth $h=td$ is $\approx 4(N+2h+2)/2^{h}$, which is $\approx 1$ at the first
admissible $t$ and falls by a factor $2^{d}$ at each subsequent $t$. The
probability that a given $(d,N)$ fails at every $t$ is therefore a product
$\prod_{j\ge0}\min(1,\,c\,2^{-jd})$, super-exponentially small, and the
expected number of failing pairs $(d,N)$ over all $N$ converges. At $d=1$,
$N=300$ the product evaluates to about $10^{-5}$, and the observed value is a
success at the very first admissible depth. The model therefore predicts
Route 4 holds with enormous margin --- which is precisely why no finite
computation will ever be evidence for it. \ev{Open} (heuristic).

\subsection{Route 5: the rationality-side rank bound remains open}

\paragraph{The exact statement that was wanted.}
\[
  \exists C\ \forall c:\N\to\N\ \bigl(c(n)\le n,\ \neg\,\mathrm{Irrational}(X_c)\bigr)
  \ \forall v>0\ \forall u\ \mathrm{IsTemperedBinaryOrbit}(c,v,u)\ \forall e:\quad
  \operatorname{rk}_e(u)\le C ,
\]
where $X_c=\sum_{n\ge1}c(n)/2^{n}$ and $\operatorname{rk}_e(u)$ is the
dimension of the span of the dyadic sections of $u$ through level $e$.
A suitable subexponential upper bound would also contradict the landed
$\varphi$-specific floor
\lean{Erdos249257.\allowbreak not_irrational_totientSeries_\allowbreak implies_unbounded_carryRank_unconditional}{TotientCarryKernelRigidity.lean:300}
\quad \ev{Lean}.  No such rationality-side upper bound is proved.

\paragraph{The right vocabulary, and the audited gap.}
Finite-dimensional span of all dyadic sections is the definition of a
$2$-regular sequence in the sense of Allouche and Shallit \ev{Cited}.  The
unconditional theorem
\lean{Erdos249257.\allowbreak not_finiteDimensional_\allowbreak span_fullTotientKernel}{TotientMahlerDefect.lean:1145}
says that $\varphi$ is not $2$-regular.  An attempted generic counterexample
chose a coefficient sequence $c(n)\le n$ with rational binary series and
non-$2$-regular dyadic kernel.  Those facts do not refute the desired bound
for its carry orbit.

The generic recurrence
\[
  v\,c(N+1)=2u(N)-u(N+1)
\]
does express each \emph{positive-residue} level-$j$ section of $c$ as a
linear combination of two level-$j$ sections of $u$.  It does not control the
zero-residue section at each level.  Across levels $1,\ldots,e$ those omitted
sections can contribute up to $e$ new directions, not a fixed $O(1)$ error.
Consequently non-$2$-regularity of $c$ alone does not imply unbounded
$\operatorname{rk}_e(u)$, and the previously claimed bound
$\operatorname{rk}_e(u)\ge2^{e-1}-3$ has no valid proof here.  The proposed
greedy set construction therefore supplies neither a counterexample nor a
reason to retire the route. \ev{Gap} \scale{uniform}
\coord{other:carry-kernel}.

The Lean theorem for $\varphi$ remains valid: it uses the special structure
and independently proved linear independence of the totient kernel, not the
invalid generic $O(1)$ bridge.  The honest status is thus asymmetric.  The
large $\varphi$-specific rank floor is proved; a rationality-driven ceiling is
open; and the generic countermodel attempt is inconclusive.

\subsection{What the shape of the wall suggests about the object}\label{sub:shape}

Everything in this subsection is inference from the assembled evidence, not
theorem. It is written because the assembly makes one inference available
that no individual reformulation does, and because refusing to draw it would
be a different kind of dishonesty than drawing it carelessly.

\paragraph{First, a fact rather than an inference: the routes are one route
in five coordinates.}
Lemma~\ref{lem:orbit} collapses the coordinate spread that made the corpus
look like many independent attacks.

\begin{center}
\begin{tabular}{p{0.16\linewidth} p{0.44\linewidth} p{0.32\linewidth}}
\hline
Route & What it asks of the binary expansion & Positions interrogated \\
\hline
1, 2 & at least $11\%$ of positions carry a digit change; mean run length
$\le 9.1$ in one block per scale & a full dyadic block of $\alpha_h$ \\
3 & one run of length $\le 3$ & positions $p-h$, $p$ prime, in $\alpha_h$ \\
4 & the tails at gap $h$ differ by more than $2(N+2h+2)2^{-h}$ & every
basepoint, some multiple of every ray, in $S$ \\
5 & an open rationality-side upper bound on the carry-orbit kernel rank &
all dyadic sections through level $e$ \\
\hline
\end{tabular}
\end{center}

\noindent
This is the honest content of the operator's ``hundred measurements of one
wall''. The measurements are of one thing, and the thing is the binary digit
sequence of $S$ and of its multiples $(2^{h}-1)S$. The routes differ in which
positions they interrogate and how uniform a margin they demand; they do not
differ in subject.

\paragraph{Second: the wall is a wall about explicit constants, not about
$\varphi$.}
Of the seven barriers, six --- B1, B2, B3, B5, B6, B7 --- are statements that
a class of proof cannot see the digits: bounded inspection, retargeting,
strengthening, bounded state, finite linear compression, coarse invariants.
None asserts that the digits misbehave. Only B4 is internal to the object,
and even B4 says that one engineering scheme self-cancels, not that the
quantity being engineered is unattainable. Meanwhile the two object-level
witnesses the corpus can produce --- the $\gamma$-splice of B1, the parity
coboundary of B7 --- and the one produced here
(Theorem~\ref{thm:lacunary}) are all lacunary or eventually periodic:
constructed, measure-zero, and highly structured. In two centuries nobody has
exhibited a naturally occurring constant with Liouville-type digit behaviour;
every known example is built to have it. So the shape of the barrier set
points at our technology, not at $S$.

\paragraph{Third: is the object behind the measurements simple or complex?}
The evidence available leans one way, and I will say how far it leans.

\emph{For ``generic, and the difficulty is entirely ours''.} (a) Every
statistic measurable is indistinguishable from a fair-coin digit model:
$\rho_h(X)\approx0.50$ against a required $0.11$; mean run length
$1.96$--$2.03$ against a permitted $9.1$; longest run $12$--$14$ over $5800$
positions against a model prediction of $\approx12.5$; block sums of order
$X^{-1/2}$ against a permitted $0.89$ (Table~\ref{tab:blocks}, \ev{Cert}). (b) The
requirements are weaker than the apparent truth by a factor $\sqrt X$; a
proof does not need to understand the digits, only to exclude a pathology.
(c) All countermodels are constructed lacunary objects. (d) The one
coordinate in which $S$ has classical structure --- the M\"obius--Mersenne
form \eqref{eq:mobmers} --- is an alternating sum of rational functions of
$2^{d}$ with no functional equation, no modularity, no continued-fraction
structure, and no algebraicity: there is no hidden object to find, which is
consistent with genericity.

\emph{Against, or at least complicating.} (a) B4's self-cancellation is
exact, unconditional, and holds at every depth $K$: the cost of forcing a
depth-$K$ residue by a Dirichlet prime is $p\ge1+2^{K-1}$ while the payoff is
amplitude $2^{K-1}$, so the trade is precisely null. An exact null of that
form is a structural fact and not an accident of parametrisation, and it is
the one place where the object itself pushes back. (b) The exact census
records residues approaching the forbidden edge --- a closest central margin
of about $2.2\cdot10^{-4}$ of the modulus at $t=100$ --- which is what a
uniform model predicts and which rules out, permanently, any argument that
proceeds by a crude uniform margin. (c) Proposition~\ref{prop:dickman} shows
that even the pieces of the frontier that are provable become true only past
$X\approx10^{30}$--$10^{40}$; a genuinely simple object would not usually
require asymptotics that begin so late.

\emph{Verdict, flagged as inference.} The evidence supports ``the digits of
$S$ are generic and the wall is technological'' over ``the object is
subtle'', but it does not determine it, and the reason it cannot is
structural: every test performed is a test a generic object passes, so
passing them is weak evidence, and the only tests that would discriminate are
exactly the ones no technique can run. What the evidence \emph{does} determine
is narrower and firmer: the difficulty of \#249 is not located in $\varphi$'s
irregularity --- $\varphi$ is the smoothest interesting multiplicative
function, with multiplicative fluctuation $\log\log Y$ where $d(n)$ has
$Y^{c/\log\log Y}$ --- but in the fact that we possess no method whatsoever
for lower-bounding digit changes of a constant that was not designed to have
them. That is why the Erd\H{o}s--Borwein constant fell in 1948 and this one
has not: $d(n)$ has spikes to deposit into the digit stream and $\varphi(n)$
does not, and in the M\"obius coordinate the signs cancel.

\paragraph{What would actually change the picture.}
One theorem, of any strength, giving a nontrivial digit statistic for $S$ or
for some $\alpha_h$: an upper bound $o(N)$ on the longest binary run in the
first $N$ positions would not by itself decide \#249, but it would be the
first evidence that the object is legible at all, and by
Corollary~\ref{cor:digitform} any bound strong enough to give a positive
proportion of digit changes in one block per scale would decide it. Failing
that, the two concrete openings identified above are: the logarithmically
averaged form of the block gap, where the entropy-decrement method is at
least in the right regime (\S\ref{sub:pivot}); and the canonical prime fibre
\eqref{eq:primefibre}, where every budget except one decorrelation estimate
has been discharged. Neither is close. Both are well posed, and both can be
started on today.
